%-----------------------------------------------------------------------------
%
%               Template for sigplanconf LaTeX Class
%
% Name:         sigplanconf-template.tex
%
% Purpose:      A template for sigplanconf.cls, which is a LaTeX 2e class
%               file for SIGPLAN conference proceedings.
%
% Guide:        Refer to "Author's Guide to the ACM SIGPLAN Class,"
%               sigplanconf-guide.pdf
%
% Author:       Paul C. Anagnostopoulos
%               Windfall Software
%               978 371-2316
%               paul@windfall.com
%
% Created:      15 February 2005
%
%-----------------------------------------------------------------------------


\documentclass[numbers]{sigplanconf}

% Uncomment when preprint off
% \documentclass[numbers]{sigplanconf}

% The following \documentclass options may be useful:

% preprint       Remove this option only once the paper is in final form.
%  9pt           Set paper in  9-point type (instead of default 10-point)
% 11pt           Set paper in 11-point type (instead of default 10-point).
% numbers        Produce numeric citations with natbib (instead of default author/year).
% authorversion  Prepare an author version, with appropriate copyright-space text.

\usepackage{amsmath}
\usepackage{amsmath}
\usepackage{amsthm}
\usepackage{amssymb}
\usepackage{multicol}
\setcounter{secnumdepth}{3}
\usepackage{xurl}
\usepackage[colorlinks=true, linkcolor=blue, linktoc=all, citecolor=black]{hyperref}
\usepackage[normalem]{ulem}
\usepackage{enumitem}
\usepackage{import}

% Better graphics handling
\usepackage{graphicx}

% Sub figures
\usepackage{subcaption}

% \setlength{\parindent}{1em}
% \setlength{\parskip}{0pt}
\usepackage{indentfirst}

% Underlining that breaks across lines
\usepackage[normalem]{ulem}

% For checkmarks
\usepackage{amssymb}

% For tables
\usepackage{booktabs}

% For multi-row cells in tables
\usepackage{multirow}

\newcommand{\cL}{{\cal L}}

\begin{document}

\special{papersize=8.5in,11in}
\setlength{\pdfpageheight}{\paperheight}
\setlength{\pdfpagewidth}{\paperwidth}

% For compatibility with auto-generated ACM eRights management
% instructions, the following alternate commands are also supported.
%\CopyrightYear{2016}
%\conferenceinfo{CONF'yy,}{Month d--d, 20yy, City, ST, Country}
%\isbn{978-1-nnnn-nnnn-n/yy/mm}\acmPrice{\$15.00}
%\doi{http://dx.doi.org/10.1145/nnnnnnn.nnnnnnn}

% Uncomment the publication rights used.
%\setcopyright{acmcopyright}
\setcopyright{acmlicensed}  % default
%\setcopyright{rightsretained}

\titlebanner{Project Check-In Preprint}        % These are ignored unless
\preprintfooter{This paper provides a structured perspective on how generative AI reshapes developer problem-solving and suggests directions for future research.}   % 'preprint' option specified.

\title{Exploration \& Acceleration: How Generative AI Reshapes Development Practices}
% \subtitle{Subtitle Text, if any\titlenote{with optional subtitle note}}

\authorinfo{Nikhil Anand}
{University of Massachusetts, Amherst}
{nikhilanand@umass.edu}
\maketitle

\begin{abstract}
    Generative AI tools such as GitHub Copilot and ChatGPT are increasingly integrated into software development workflows, yet their influence on developer problem-solving strategies remains insufficiently understood. This paper presents a literature survey of recent studies examining how AI Assisted programming reshapes developer behavior across different stages of the Software Development Lifecycle (SDLC). We first propose a taxonomy of developer AI interaction paradigms, including suggestion-based, conversational, generative, and retrieval-oriented assistance. Using this taxonomy, we analyze the impact of Generative AI on requirements elicitation, coding, debugging, and documentation workflows. The surveyed literature shows that AI systems significantly improve productivity and implementation speed while shifting developers from direct code authorship toward solution evaluation and refinement. However, these benefits are accompanied by limitations related to reliability, overreliance, shallow reasoning, and reduced code comprehension. Overall, the findings suggest that the future of software development lies in effective Human AI collaboration rather than fully autonomous software generation.
\end{abstract}

% 2012 ACM Computing Classification System (CSS) concepts
% Generate at 'http://dl.acm.org/ccs/ccs.cfm'.
\begin{CCSXML}
    <ccs2012>
    <concept>
    <concept_id>10011007.10011074.10011081</concept_id>
    <concept_desc>Software and its engineering~Software development process management</concept_desc>
    <concept_significance>500</concept_significance>
    </concept>
    <concept>
    <concept_id>10010147.10010178.10010219.10010220</concept_id>
    <concept_desc>Computing methodologies~Multi-agent systems</concept_desc>
    <concept_significance>300</concept_significance>
    </concept>
    <concept>
    <concept_id>10010147.10010178.10010179.10003352</concept_id>
    <concept_desc>Computing methodologies~Information extraction</concept_desc>
    <concept_significance>100</concept_significance>
    </concept>
    <concept>
    <concept_id>10010147.10010178.10010179.10010182</concept_id>
    <concept_desc>Computing methodologies~Natural language generation</concept_desc>
    <concept_significance>100</concept_significance>
    </concept>
    </ccs2012>
\end{CCSXML}

\ccsdesc[500]{Software and its engineering~Software development process management}
\ccsdesc[300]{Computing methodologies~Multi-agent systems}
\ccsdesc[100]{Computing methodologies~Information extraction}
\ccsdesc[100]{Computing methodologies~Natural language generation}
% end generated code

% Legacy 1998 ACM Computing Classification System categories are also
% supported, but not recommended.
%\category{CR-number}{subcategory}{third-level}[fourth-level]
%\category{D.3.0}{Programming Languages}{General}
%\category{F.3.2}{Logics and Meanings of Programs}{Semantics of Programming Languages}[Program analysis]

% \keywords
% keyword1, keyword2

\import{content/}{introduction.tex}
\import{content/}{taxonomy.tex}


% The 'abbrvnat' bibliography style is recommended.

\bibliographystyle{abbrvnat}
\bibliography{refs.bib}

% Paste babel file here


\end{document}
